\Legend
小明有一个长度为 $N$ 的字符串 $S$，其中每个字符都属于 $\{\texttt{a},\texttt{b},\texttt{c},\texttt{d},\texttt{e}\}$。

给定字符串 $T$，长度为 $M$，定义其哈希值为

$$\sum_{i=1}^{M}(T_i-\texttt{a}+1)\cdot1331^{M-i}\pmod{2^{32}},$$

其中字符之差按 ASCII 码计算。

对于每个询问 $l,r,k$，考虑区间 $S_lS_{l+1}\cdots S_r$ 的所有\textbf{不同}子序列。这些子序列必须以字符 $S_l$ 开头，以字符 $S_r$ 结尾，且长度至少为 $2$。不同只按生成的字符串内容判断，与选取下标的方式无关。把合法字符串按字典序从小到大排序，输出第 $k$ 个字符串的哈希值；如果合法字符串不足 $k$ 个，输出 $-1$。

子序列可由原区间删除任意若干字符得到，未删除字符的顺序保持不变。字典序采用通常的规则：在首个不同位置，字符较小者更小；若一个字符串是另一个的前缀，较短者更小。

\Input

第一行输入整数 $N,Q$。第二行输入长度为 $N$ 的字符串 $S$。接下来 $Q$ 行，每行输入三个整数 $l,r,k$。

$1\le N,Q\le10^5$，$1\le l\le r\le N$，$1\le k\le10^9$；$S$ 只含字符 \tt{a} 到 \tt{e}。

\Output

每个询问输出一行，表示答案。

\Examples

\begin{example}
\exmp{2 3
aa
1 2 1
1 2 2
1 1 1}{1332
-1
-1}
\end{example}

\begin{example}
\exmp{10 10
ddcbbbbbbc
2 7 2
1 9 5
1 8 7
1 9 2
3 8 2
2 8 1
1 7 2
1 10 5
1 8 1
1 10 10}{7088908
272200126
847173519
7088908
5317347
5326
7088908
845401959
5326
3153146560}
\end{example}
